%% Đề 42 -- de-tuyen-sinh-lop-10-thpt-mon-toan-nam-2014-2015-so-gddt-lang-son.pdf
\deheader{ĐỀ SỐ 42}%
  {Kỳ thi Tuyển sinh vào lớp 10 THPT — Năm học 2014--2015}%
  {Sở GD\&ĐT Lạng Sơn | Chỉ phần Đại số}
\addcontentsline{toc}{section}{Đề 42 — Tuyển sinh vào lớp 10 THPT — Năm học 2014--2015}

\begin{de_cau}{Câu 1}{2,0 điểm}
\begin{enumerate}[label=\alph*)]
  \item Tính: $A = 2\sqrt{3}+\sqrt{12}-\sqrt{27}$.
  \item Rút gọn: $B = \dfrac{\sqrt{x}+1}{\sqrt{x}-1}-\dfrac{2}{\sqrt{x}+1}-\dfrac{2}{x-1}$ với $x\ge0$, $x\neq1$.
\end{enumerate}
\end{de_cau}

\begin{loi_giai_cau}
\textbf{a)} $A=2\sqrt{3}+2\sqrt{3}-3\sqrt{3}=\boxed{\sqrt{3}}$.

\textbf{b)} Quy đồng mẫu $(\sqrt{x}-1)(\sqrt{x}+1)$:

$B=\dfrac{(\sqrt{x}+1)^2-2(\sqrt{x}-1)-2}{(\sqrt{x}-1)(\sqrt{x}+1)}=\dfrac{x+2\sqrt{x}+1-2\sqrt{x}+2-2}{x-1}=\dfrac{x+1}{x-1}$.

Vậy $B=\boxed{\dfrac{x+1}{x-1}}$.
\end{loi_giai_cau}

\begin{de_cau}{Câu 2}{2,0 điểm}
\begin{enumerate}[label=\alph*)]
  \item Giải phương trình: $x^2-5x+6=0$.
  \item Giải hệ phương trình:
  \[\begin{cases} x-2y=-1 \\ 2x+y=7 \end{cases}.\]
\end{enumerate}
\end{de_cau}

\begin{loi_giai_cau}
\textbf{a)} $(x-2)(x-3)=0\Rightarrow\boxed{x=2}$ hoặc $\boxed{x=3}$.

\textbf{b)} Nhân pt(1) với 2: $2x-4y=-2$. Trừ pt(2): $-5y=-9\Rightarrow y=\dfrac{9}{5}$... Hoặc: $x=2y-1$; $2(2y-1)+y=7\Rightarrow5y=9\Rightarrow y=\frac{9}{5}$, $x=\frac{13}{5}$.

Vậy $\boxed{(x;y)=\left(\dfrac{13}{5};\dfrac{9}{5}\right)}$.
\end{loi_giai_cau}

\begin{de_cau}{Câu 3}{2,0 điểm}
Cho phương trình $x^2-2(m-1)x+m^2-4=0$ (tham số $m$).
\begin{enumerate}[label=\alph*)]
  \item Tìm $m$ để phương trình có nghiệm kép. Tính nghiệm kép đó.
  \item Tìm $m$ để phương trình có hai nghiệm $x_1,x_2$ với $x_1+x_2=2x_1x_2$.
\end{enumerate}
\end{de_cau}

\begin{loi_giai_cau}
\textbf{a)} $\Delta'=(m-1)^2-(m^2-4)=-2m+5=0\Rightarrow m=\dfrac{5}{2}$.

Nghiệm kép: $x=m-1=\boxed{\dfrac{3}{2}}$.

\textbf{b)} $x_1+x_2=2(m-1)$, $x_1x_2=m^2-4$.

$2(m-1)=2(m^2-4)\Rightarrow m-1=m^2-4\Rightarrow m^2-m-3=0\Rightarrow m=\dfrac{1\pm\sqrt{13}}{2}$.

Cần $\Delta'\ge0$: $m\le\dfrac{5}{2}$. Cả hai nghiệm $\le\dfrac{5}{2}$. Vậy $\boxed{m=\dfrac{1\pm\sqrt{13}}{2}}$.
\end{loi_giai_cau}

\begin{de_cau}{Câu 5}{0,5 điểm}
Cho $x, y > 0$ và $x + y = 1$. Tìm giá trị lớn nhất của $S = x^2y + xy^2$.
\end{de_cau}

\begin{loi_giai_cau}
$S=xy(x+y)=xy\le\dfrac{(x+y)^2}{4}=\dfrac{1}{4}$.

Dấu $=$ khi $x=y=\dfrac{1}{2}$. Vậy $S_{\max}=\boxed{\dfrac{1}{4}}$.
\end{loi_giai_cau}
